\documentclass[10pt, aspectratio=169]{beamer}

\usepackage[utf8]{inputenc}
\usepackage[T1]{fontenc}
\usepackage[english]{babel}
\usepackage{lmodern}

\usepackage{amsmath, amssymb}
\usepackage{booktabs}
\usepackage{graphicx} % Aggiunto per inserire i grafici
\graphicspath{{immagini/}} % Cartella in cui cercare le immagini
\usepackage{xcolor}
\usepackage{pgfplots}
\pgfplotsset{compat=1.18}

% Palette
\definecolor{NavyPrimary}{RGB}{18, 43, 73}
\definecolor{TealAccent}{RGB}{0, 128, 128}
\definecolor{SlateNeutral}{RGB}{71, 85, 105}
\definecolor{LightTint}{RGB}{241, 245, 249}
\definecolor{PlotBlue}{RGB}{30, 136, 229}
\definecolor{PlotRed}{RGB}{229, 57, 53}
\definecolor{PlotOrange}{RGB}{251, 140, 0}
\definecolor{QuestionRed}{RGB}{190, 30, 45} % Colore rosso per le domande

% Comando per le domande in rosso
\newcommand{\question}[1]{{\color{QuestionRed}\textbf{#1}}}

% Clean Beamer styling
\usetheme{default}
\usecolortheme{default}
\setbeamertemplate{navigation symbols}{}
\setbeamercolor{title}{fg=NavyPrimary}
\setbeamercolor{subtitle}{fg=TealAccent}
\setbeamercolor{author}{fg=SlateNeutral}
\setbeamercolor{date}{fg=SlateNeutral}
\setbeamercolor{frametitle}{fg=NavyPrimary, bg=LightTint}
\setbeamercolor{structure}{fg=NavyPrimary}

\setbeamerfont{title}{size=\LARGE, series=\bfseries}
\setbeamerfont{subtitle}{size=\large}
\setbeamerfont{frametitle}{size=\normalsize, series=\bfseries}

\setbeamertemplate{footline}{
	\leavevmode%
	\hbox{%
		\begin{beamercolorbox}[wd=.75\paperwidth,ht=2.5ex,dp=1ex,leftskip=1.5em]{date}%
			\usebeamerfont{date}ASSIGNMENT 1 -- Communications Systems
		\end{beamercolorbox}%
		\begin{beamercolorbox}[wd=.25\paperwidth,ht=2.5ex,dp=1ex,rightskip=1.5em]{date}%
			\raggedleft\usebeamerfont{date}\insertframenumber{} / \inserttotalframenumber
		\end{beamercolorbox}%
	}%
	\vskip2pt%
}

\title{ASSIGNMENT 2}
\subtitle{Communications Systems}
\author{Alessandro Felicetti 375462 \and Ivan Pizzulli 368369 \\ Mattia Cilio 377785 \and Luca Bortoluzzi 378057}
\date{}

\begin{document}
	
	% Slide 1
	\begin{frame}[plain]
		\vspace{1.0cm}
		\begin{flushleft}
			{\color{TealAccent}\rule{2.5cm}{2pt}}\\[0.5cm]
			{\usebeamerfont{title}\usebeamercolor[fg]{title}\inserttitle}\\[0.2cm]
			{\usebeamerfont{subtitle}\usebeamercolor[fg]{subtitle}\insertsubtitle}\\[1.0cm]
			{\usebeamerfont{author}\usebeamercolor[fg]{author}\insertauthor}
		\end{flushleft}
	\end{frame}
	
	% Slide 2
	\begin{frame}{Exercises}
		\begin{itemize}
			\item Gold codes
			\item CDMA
			\item Link Budget
			\item ...
			\item ...
		\end{itemize}
	\end{frame}
	
	% Slide 3
	\begin{frame}[plain]
		\vfill
		\centering
		{\color{TealAccent}\rule{2cm}{2pt}}\\[0.4cm]
		{\LARGE\bfseries\color{NavyPrimary} Exercise 1}\\[0.2cm]
		{\large\color{SlateNeutral} Gold codes}
		\vfill
	\end{frame}
	
	% Slide 4
	\begin{frame}[plain]
		\vfill
		\centering
		{\color{TealAccent}\rule{2cm}{2pt}}\\[0.4cm]
		{\LARGE\bfseries\color{NavyPrimary} Exercise 1 - Part 1}\\[0.2cm]
		\vfill
	\end{frame}
	
	% Slide 5
	\begin{frame}{Exercise 1 - Part 1}
		\question{1. Generate a Gold sequence with index i $ \geq 0$}
		
		\vspace{0.1cm}
		\begin{itemize}
			\item Using the \texttt{comm.GoldSequence} object in MATLAB with $m = 7$ shift register stages (sequence length $L = 2^7 - 1 = 127$), we generated a Gold code setting \texttt{Index} $= 1 \ge 0$.
			\item The generator uses the preferred pair of primitive polynomials $x^7 + x^4 + 1$ and $x^7 + x + 1$, both initialized with state $[0\,0\,0\,0\,0\,0\,1]$, and the resulting binary sequence is mapped to bipolar form $\{+1, -1\}$ via $1 - 2x$.
		\end{itemize}
		
	\end{frame}
	
	% Slide 6
	\begin{frame}{Exercise 1 - Part 1}
		\question{2. Compute and plot the cyclic auto-correlation}
		
		\vspace{0.1cm}
		\begin{itemize}
			\item Autocorrelation function for Gold sequence $(i=1)$
		\end{itemize}
		
		\vspace{0.2cm}
		\begin{center}
			\includegraphics[width=0.72\textwidth, height=0.58\textheight, keepaspectratio]{A2_1_2.png}
		\end{center}
		
	\end{frame}
	
	% Slide 7
	\begin{frame}{Exercise 1 - Part 1}
		\question{3. Repeat for $i = -1$ and $i = -2$}
		
		\vspace{0.2cm}
		\begin{columns}[T]
			\begin{column}{0.48\textwidth}
				\centering
				\textbf{Index $i = -1$}\\[0.15cm]
				\includegraphics[width=\textwidth, height=0.55\textheight, keepaspectratio]{A2_1_3_m1.png}
			\end{column}
			\begin{column}{0.48\textwidth}
				\centering
				\textbf{Index $i = -2$}\\[0.15cm]
				\includegraphics[width=\textwidth, height=0.55\textheight, keepaspectratio]{A2_1_3_m2.png}
			\end{column}
		\end{columns}
		
	\end{frame}
	
	% Slide 8
	\begin{frame}{Exercise 1 - Part 1}
		\question{4. Comment on the results}
	
		\vspace{0.1cm}
		\begin{itemize}
			\item \textbf{$i = 1$ (Gold sequence):} The autocorrelation peaks at $L = 127$ ($\tau = 0$) and takes the 3 theoretical values $\{-17, -1, 15\}$ elsewhere, matching $t(m) = 2^{(m+1)/2} + 1 = 17$ for $m = 7$.
			\item \textbf{$i = -1$ and $i = -2$ (m-sequences):} MATLAB outputs the two generating m-sequences. Even though their bits are different, they share the exact same ideal two-level autocorrelation ($127$ at $\tau = 0$, $-1$ everywhere else).
			\item Pure m-sequences have ideal autocorrelation for synchronization, but give only 2 codes. Gold codes trade slightly worse autocorrelation sidelobes for 129 low-interference codes, making them ideal for multi-user CDMA.
		\end{itemize}
	
	\end{frame}
	
	% Slide 9
	\begin{frame}[plain]
		\vfill
		\centering
		{\color{TealAccent}\rule{2cm}{2pt}}\\[0.4cm]
		{\LARGE\bfseries\color{NavyPrimary} Exercise 1 - Part 2}\\[0.2cm]
		\vfill
	\end{frame}
	
	% Slide 10
	\begin{frame}{Exercise 1 - Part 2}
		\question{6. Compute and plot their cyclic cross-correlation}
		
		\vspace{0.1cm}
		\begin{itemize}
			\item Cross correlation of two Gold sequence with random $i \geq 0$.
		\end{itemize}
		
		\vspace{0.2cm}
		\begin{center}
			\includegraphics[width=0.72\textwidth, height=0.58\textheight, keepaspectratio]{A2_1_6.png}
		\end{center}
		
	\end{frame}
	
	
	% Slide 11
	\begin{frame}{Exercise 1 - Part 2}
		\question{8. Compute and plot their cyclic cross-correlation}
		
		\vspace{0.1cm}
		\begin{itemize}
			\item Two different m-sequences with the same length using two distinct primitive polynomials that do not form a preferred pair.
		\end{itemize}
		
		\vspace{0.2cm}
		\begin{center}
			\includegraphics[width=0.72\textwidth, height=0.58\textheight, keepaspectratio]{A2_1_8.png}
		\end{center}
		
	\end{frame}
	
	% Slide 12
	\begin{frame}{Exercise 1 - Part 2}
		\question{9. Compare and discuss the two results}
		
		\vspace{0.1cm}
		\begin{itemize}
			\item While both cross-correlation waveforms appear noise-like across the lags, the two Gold sequences exhibit a well-bounded three-level structure taking only the values $\{-17, -1, 15\}$ (corresponding to $\{-t, -1, t-2\}$ for $m = 7$).
			\item In contrast, the cross-correlation between the two non-preferred m-sequences ($x^7 + x^4 + 1$ and $x^7 + x^2 + 1$) loses this three-level property and fluctuates across many discrete levels, reaching significantly larger values (up to $+25$ and $-23$, well beyond the Gold code bounds)[cite: 2].
			\item Because these two primitive polynomials do not form a preferred pair, their maximum cross-correlation is no longer bounded by $t(m) = 17$, resulting in larger interference peaks and worse multi-user separation.
		\end{itemize}
		
	\end{frame}
	
	% Slide 13
	\begin{frame}[plain]
		\vfill
		\centering
		{\color{TealAccent}\rule{2cm}{2pt}}\\[0.4cm]
		{\LARGE\bfseries\color{NavyPrimary} Exercise 1 - Part 3}\\[0.2cm]
		\vfill
	\end{frame}
	
	% Slide 14
	\begin{frame}{Exercise 1 - Part 3}
		\question{10. Plot the three versions of the Welch bound and the Sidelnikov bound for $N = 127$ and $1 \leq K \leq 200$}
		
		\vspace{0.2cm}
		\begin{center}
			\includegraphics[width=0.72\textwidth, height=0.58\textheight, keepaspectratio]{A2_1_10.png}
		\end{center}
		
	\end{frame}
		
	% Slide 15
	\begin{frame}{Exercise 1 - Part 3}
		\question{11. Consider the entire set of $K = 129$ Gold sequences and compute $r_M$ (explain how you computed its value).}
		
		\vspace{0.1cm}
		\begin{itemize}
			\item We generated the entire set of $K = 129$ Gold sequences of length $N = 127$ by varying the \texttt{Index} parameter from $-2$ to $126$ ($N - 1$), including the two base m-sequences ($-2$ and $-1$) from $x^7 + x^4 + 1$ and $x^7 + x + 1$, along with their $127$ cyclic-shifted combinations ($0$ to $126$).
			\item After converting the sequences to bipolar form ($+1, -1$), we computed the periodic autocorrelation for all non-zero lags (excluding the peak of $N = 127$ at lag $0$) and the periodic cross-correlation across all sequence pairs and lags.
			\item Taking the maximum absolute value across all these correlations gave $r_M = 17$, which we then plotted as the point $(129, 17)$ on the bound figure.
		\end{itemize}
		
	\end{frame}
	
	% Slide 16
	\begin{frame}{Exercise 1 - Part 3}
		\question{12. On the bound figure, add a point corresponding to $r_M$.}
		
		\vspace{0.2cm}
		\begin{center}
			\includegraphics[width=0.72\textwidth, height=0.58\textheight, keepaspectratio]{A2_1_12.png}
		\end{center}
		
	\end{frame}
	
	% Slide 17
	\begin{frame}{Exercise 1 - Part 3}
		\question{13. Comment on the result.}
		
		\vspace{0.1cm}
		\begin{itemize}
			\item For $K \le 127$ ($K \le N$), the general Welch bound ($s \le 10$) and the simplified version ($s = 1$) coincide and approach $\lceil\sqrt{N}\rceil = 12$. For $K > N$, the $s = 1$ bound saturates at $12$, while the general bound keeps increasing as higher-order terms ($s \ge 2$) become dominant.
			\item The Sidelnikov bound stays strictly above all three Welch curves for every $K$, since it gives a tighter lower bound by accounting for the binary bipolar nature of the sequences, growing in discrete steps due to the ceiling operations and maximization over $s$.
			\item At $K = 129$, the Gold sequence set achieves $r_M = 17$, sitting just $1$ unit above the Sidelnikov bound ($16$) and above the Welch bounds ($13$ and $12$), showing that Gold codes are near-optimal while still providing a large family ($K = N + 2$) with bounded three-valued correlation.
		\end{itemize}
		
	\end{frame}
	
	% Slide 18
	\begin{frame}[plain]
		\vfill
		\centering
		{\color{TealAccent}\rule{2cm}{2pt}}\\[0.4cm]
		{\LARGE\bfseries\color{NavyPrimary} Exercise 1 - Part 4}\\[0.2cm]
		\vfill
	\end{frame}
	
	% Slide 19
	\begin{frame}{Exercise 1 - Part 4}
		\question{14. Starting from the Gold set of $K = 129$ sequences, remove the last chip of each sequence to obtain length $N = 126$. Compute $r_M$ using cyclic correlations of period $126$ (explain how you computed its value).}
		
		\vspace{0.1cm}
		\begin{itemize}
			\item Starting from the full bipolar Gold set of $K = 129$ sequences of length $127$, we truncated each sequence by removing its last chip, obtaining a new set of length $N = 126$.
			\item Using $126$-point FFT/IFFT, we computed the periodic autocorrelation for all non-zero lags (excluding the lag-$0$ peak) to find the maximum off-peak value $R_A$, and the periodic cross-correlation across all sequence pairs and lags to find $R_C$.
			\item Taking $r_M = \max(R_A, R_C)$ across the entire truncated set yielded $r_M = 26$.
		\end{itemize}
		
	\end{frame}
	
	% Slide 20
	\begin{frame}{Exercise 1 - Part 4}
		\question{15. Compute and plot the bounds for $N = 126$ and $1 \leq K \leq 200$.}
		
		\vspace{0.2cm}
		\begin{center}
			\includegraphics[width=0.72\textwidth, height=0.58\textheight, keepaspectratio]{A2_1_15.png}
		\end{center}
		
	\end{frame}
	
	% Slide 21
	\begin{frame}{Exercise 1 - Part 4}
		\question{16. Add $r_M$ to the plot.}
		
		\vspace{0.2cm}
		\begin{center}
			\includegraphics[width=0.72\textwidth, height=0.58\textheight, keepaspectratio]{A2_1_16.png}
		\end{center}
		
	\end{frame}
	
	% Slide 22
	\begin{frame}{Exercise 1 - Part 4}
		\question{17. Comment on the result.}
	
		\vspace{0.1cm}
		\begin{itemize}
			\item For $N = 126$ and $K = 129$, the maximum cyclic correlation reaches $r_M = 26$, which is significantly higher than both the binary Sidelnikov bound ($16$) and the Welch bounds ($12$).
			\item Unlike the original $N = 127$ case where Gold sequences are near-optimal thanks to the Galois field $GF(2^m)$ structure and preferred pairs removing even a single chip breaks the native $N = 2^m - 1$ periodicity and the balance between $+1$ and $-1$.
			\item As a result, the three-valued cyclic correlation property $\{-t, -1, t-2\}$ is completely lost, and out-of-phase autocorrelation and cross-correlation peaks increase sharply, making truncated sequences sub-optimal for asynchronous CDMA.
		\end{itemize}
	
	\end{frame}
	
	% Slide 23
	\begin{frame}[plain]
		\vfill
		\centering
		{\color{TealAccent}\rule{2cm}{2pt}}\\[0.4cm]
		{\LARGE\bfseries\color{NavyPrimary} Exercise 1 - Part 5 - Walsh–Hadamard Sequences}\\[0.2cm]
		\vfill
	\end{frame}
	
	% Slide 24
	\begin{frame}{Exercise 1 - Part 5}
		\question{20. Compute $r_A$, $r_C$ and $r_M$ using cyclic correlations.}
		
		\vspace{0.1cm}
		\begin{itemize}
			\item Using $128$-point FFT/IFFT across all $K = 128$ rows of the Walsh-Hadamard matrix, we computed the periodic autocorrelation (excluding the lag-$0$ peak of $N = 128$) and the periodic cross-correlation for every pair of sequences and all cyclic shifts.
			\item Taking the maximum absolute values gave $r_A = 128$, $r_C = 128$, and therefore $r_M = \max(r_A, r_C) = 128$.
			\item This happens because the all-ones first row has an autocorrelation of $128$ at every shift, and cyclically shifting certain rows aligns them completely with other rows in the matrix.
		\end{itemize}
		
	\end{frame}
	
	% Slide 25
	\begin{frame}{Exercise 1 - Part 5}
		\question{21. Compare with Welch and Sidelnikov bounds for $N = K = 128$.}
		
		\vspace{0.1cm}
		\begin{itemize}
			\item Evaluating the theoretical lower bounds at $N = K = 128$, we obtain $12$ for all Welch formulations (general Welch bound, $s = 1$ formula, and $\lceil\sqrt{N}\rceil$ approximation) and $16$ for the binary Sidelnikov bound.
			\item The Walsh-Hadamard family achieves $r_M = 128$, which is far worse than both the Sidelnikov ($16$) and Welch ($12$) bounds, confirming that these sequences are poorly suited for asynchronous periodic correlation.
		\end{itemize}
		
	\end{frame}
	
	% Slide 26
	\begin{frame}{Exercise 1 - Part 5}
		\question{22. Comment on the results. Discuss downlink and uplink use cases. Discuss synchronization issues.}
		
		\vspace{0.1cm}
		\begin{itemize}
			\item Walsh-Hadamard sequences are strictly orthogonal only at zero delay ($\tau = 0$), while for $\tau \neq 0$ their correlation reaches $r_A = r_C = 128$, far above the Sidelnikov ($16$) and Welch ($12$) bounds.
			\item They work in a synchronous single-transmitter downlink ($\tau = 0$), whereas in the uplink distributed users arrive asynchronously ($\tau \neq 0$) due to different propagation delays and local clocks.
			\item Since $r_A = r_C = 128$, any timing offset or multipath delay ($\tau \neq 0$) destroys orthogonality and prevents code acquisition (no sharp peak at $\tau = 0$), requiring a PN scrambling code in cellular systems.
			\item \textbf{GPS} is a key counterexample: although it is a downlink, satellites lie at different distances from the receiver, so signals arrive asynchronously ($\tau \neq 0$) and the receiver must estimate each delay $\tau$, requiring Gold codes instead.
		\end{itemize}
		
	\end{frame}
	
	% Slide 27
	\begin{frame}[plain]
		\vfill
		\centering
		{\color{TealAccent}\rule{2cm}{2pt}}\\[0.4cm]
		{\LARGE\bfseries\color{NavyPrimary} Exercise 2 - CDMA}\\[0.2cm]
		\vfill
	\end{frame}
	
	% Slide 28
	\begin{frame}[plain]
		\vfill
		\centering
		{\color{TealAccent}\rule{2cm}{2pt}}\\[0.4cm]
		{\LARGE\bfseries\color{NavyPrimary} Exercise 2 - Part 1}\\[0.2cm]
		\vfill
	\end{frame}
	
	% Slide 29
	\begin{frame}{Exercise 2 - Part 1}
		\question{1. Plot the analytical bit error probability $Pb(e)$ versus $Eb/N0$ (in dB), from $0$ to $12$ dB.}
		
		\vspace{0.2cm}
		\begin{center}
			\includegraphics[width=0.72\textwidth, height=0.58\textheight, keepaspectratio]{A2_2_1.png}
		\end{center}
		
	\end{frame}
	
	% Slide 30
	\begin{frame}{Exercise 2 - Part 1}
		\question{2. Plot the simulated BER together with the analytical curve on the same graph.}
		
		\vspace{0.2cm}
		\begin{center}
			\includegraphics[width=0.72\textwidth, height=0.58\textheight, keepaspectratio]{A2_2_2.png}
		\end{center}
		
	\end{frame}
	
	% Slide 31
	\begin{frame}{Exercise 2 - Part 1}
		\question{3. Comment on the results.}
		
		\vspace{0.1cm}
		\begin{itemize}
			\item The monte Carlo simulation on 100 errors (red curve) follows almost perfectly the Analytical curve (blue curve) with a little misalignment on 2dB
		\end{itemize}
		
	\end{frame}
	
	% Slide 32
	\begin{frame}[plain]
		\vfill
		\centering
		{\color{TealAccent}\rule{2cm}{2pt}}\\[0.4cm]
		{\LARGE\bfseries\color{NavyPrimary} Exercise 2 - Part 2}\\[0.2cm]
		\vfill
	\end{frame}
	
	% Slide 33
	\begin{frame}{Exercise 2 - Part 2}
		\question{4. Plot the analytical curve.}
		
		\vspace{0.2cm}
		\begin{center}
			\includegraphics[width=0.72\textwidth, height=0.58\textheight, keepaspectratio]{A2_2_4.png}
		\end{center}
		
	\end{frame}
	
	% Slide 34
	\begin{frame}{Exercise 2 - Part 2}
		\question{5. Plot the simulated curve.}
		
		\vspace{0.2cm}
		\begin{center}
			\includegraphics[width=0.72\textwidth, height=0.58\textheight, keepaspectratio]{A2_2_5.png}
		\end{center}
		
	\end{frame}
	
	% Slide 35
	\begin{frame}{Exercise 2 - Part 2}
		\question{6. Comment on the results.}
		
		\vspace{0.1cm}
		\begin{itemize}
			\item The simulated BER curve for the single-user DSSS system closely matches the analytical 2-PAM curve $P_b(e) = \frac{1}{2}\operatorname{erfc}(\sqrt{E_b/N_0})$, showing the same performance observed in Part 1.
			\item Since only one user is active with no multi-user interference, spreading each bit over $N = 12$ unit-energy chips ($E_b = N$) and despreading at the receiver corresponds to projecting the received signal onto the normalized basis vector $\underline{c}_1'/\sqrt{N}$.
			\item At the same $E_b/N_0$, spreading neither improves nor degrades the error rate over an AWGN channel; its role is instead to enable multi-user separation in CDMA.
			\item Minor fluctuations at low $E_b/N_0$ ($0\text{–}3\text{ dB}$) are solely due to the statistical variance of the Monte Carlo estimation when stopping at $100$ bit errors.
		\end{itemize}
		
	\end{frame}
	
	% Slide 36
	\begin{frame}[plain]
		\vfill
		\centering
		{\color{TealAccent}\rule{2cm}{2pt}}\\[0.4cm]
		{\LARGE\bfseries\color{NavyPrimary} Exercise 2 - Part 3}\\[0.2cm]
		\vfill
	\end{frame}
	
	% Slide 37a
	\begin{frame}{Exercise 2 - Part 3}
		\question{7. Which is the parameter $p$ that determines the performance of this scenario? (1/2)}
		
		\vspace{0.1cm}
		\begin{itemize}
			\item Over an AWGN channel with chip noise variance $\sigma^2 = N_0/2$ and unit-energy chips ($E_b = N = 12$), the received synchronous vector at user $U_1$ is:
			\[
			\underline{r} = b_1'\underline{c}_1' + b_2'\underline{c}_2' + \underline{n}
			\]
			where $b_1', b_2' \in \{-1, +1\}$ are equiprobable bits.
			\item Despreading by projecting $\underline{r}$ onto $\underline{c}_1'$ yields the scalar decision variable:
			\[
			y = \underline{r} \cdot \underline{c}_1' = b_1'(\underline{c}_1' \cdot \underline{c}_1') + b_2'(\underline{c}_2' \cdot \underline{c}_1') + (\underline{n} \cdot \underline{c}_1')
			\]
			\item Since $c'_{1,k} \in \{-1, +1\}$, we have $\underline{c}_1' \cdot \underline{c}_1' = N$, so the desired signal term is $b_1'N$ and the projected noise is $\tilde{n} = \underline{n} \cdot \underline{c}_1' \sim \mathcal{N}(0, N \cdot N_0/2)$.
		\end{itemize}
		
	\end{frame}
	
	% Slide 37b
	\begin{frame}{Exercise 2 - Part 3}
		\question{7. Which is the parameter $p$ that determines the performance of this scenario? (2/2)}
		
		\vspace{0.1cm}
		\begin{itemize}
			\item Both the desired signal ($b_1'N$) and the noise variance ($N N_0/2$) depend solely on the fixed spreading factor $N$ and $N_0$.
			\item Because $b_2' \in \{-1, +1\}$ is equiprobable, the sign of the inner product $\underline{c}_2' \cdot \underline{c}_1'$ is absorbed by $b_2'$, so the Multiple Access Interference (MAI) term $b_2'(\underline{c}_2' \cdot \underline{c}_1')$ takes values $\pm |\underline{c}_1' \cdot \underline{c}_2'|$ with equal probability $1/2$.
			\item Thus, the performance is governed by the magnitude of the zero-lag cross-correlation:
			\[
			p = |\underline{c}_1' \cdot \underline{c}_2'| = \left|\sum_{k=1}^{N} c'_{1,k} c'_{2,k}\right| \in \{0, 2, 4, \dots, N\}
			\]
			where $p = 0$ gives zero MAI (single-user 2-PAM), while increasing $p$ toward $N = 12$ degrades BER.
		\end{itemize}
		
	\end{frame}
	
	% Slide 38a
	\begin{frame}{Exercise 2 - Part 3}
		\question{8. Derive analytically the performance curve of this scenario with two users (and no cyclic shift). (1/2)}
		
		\vspace{0.1cm}
		\begin{itemize}
			\item By antipodal symmetry ($P(b_1'=+1) = P(b_1'=-1) = 1/2$), we set $b_1' = +1$, so $P_b(e) = P(y < 0 \mid b_1' = +1)$.
			\item Conditioned on $b_2' \in \{-1, +1\}$, $y = N + b_2'(\underline{c}_2' \cdot \underline{c}_1') + \tilde{n}$ is Gaussian with mean $\mu(b_2') = N + b_2'(\underline{c}_2' \cdot \underline{c}_1')$ and standard deviation $\sigma_{\tilde{n}} = \sqrt{N \cdot N_0/2}$, giving:
			\[
			P_b(e \mid b_2') = Q\!\left(\frac{\mu(b_2')}{\sigma_{\tilde{n}}}\right) = \frac{1}{2}\operatorname{erfc}\!\left(\frac{N + b_2'(\underline{c}_2' \cdot \underline{c}_1')}{\sqrt{N \cdot N_0}}\right)
			\]
			\item Substituting $E_b = N$ ($\sqrt{N \cdot N_0} = \sqrt{E_b N_0}$) and factoring out $N$ yields:
			\[
			P_b(e \mid b_2') = \frac{1}{2}\operatorname{erfc}\!\left(\sqrt{\frac{E_b}{N_0}}\left[1 + b_2'\frac{\underline{c}_2' \cdot \underline{c}_1'}{N}\right]\right)
			\]
		\end{itemize}
		
	\end{frame}
	
	% Slide 38b
	\begin{frame}{Exercise 2 - Part 3}
		\question{8. Derive analytically the performance curve of this scenario with two users (and no cyclic shift). (2/2)}
		
		\vspace{0.1cm}
		\begin{itemize}
			\item Averaging over the two equiprobable values of $b_2'$ via the Law of Total Probability:
			\[
			P_b(e) = \frac{1}{2}P_b(e \mid b_2' = +1) + \frac{1}{2}P_b(e \mid b_2' = -1)
			\]
			\item Regardless of the sign of $\underline{c}_2' \cdot \underline{c}_1'$, one bit value produces constructive interference ($+p$) and the other produces destructive interference ($-p$), where $p = |\underline{c}_1' \cdot \underline{c}_2'|$:
			\[
			P_b(e) = \frac{1}{4}\operatorname{erfc}\!\left(\sqrt{\frac{E_b}{N_0}}\left(1 + \frac{p}{N}\right)\right) + \frac{1}{4}\operatorname{erfc}\!\left(\sqrt{\frac{E_b}{N_0}}\left(1 - \frac{p}{N}\right)\right)
			\]
			\item For $p = 0$ we recover 2-PAM; for $p = 2, 4$ the SNR penalty is $1.58\text{ dB}$ and $3.52\text{ dB}$; for $p = 12$ destructive interference cancels the signal, creating an error floor at $P_b(e) = 0.25$.
		\end{itemize}
		
	\end{frame}
	
	% Slide 39a
	\begin{frame}{Exercise 2 - Part 3}
		\question{9. Derive an approximate performance curve based on $p$. (1/2)}
		
		\vspace{0.1cm}
		\begin{itemize}
			\item At moderate-to-high $E_b/N_0$, the complementary error function decays super-exponentially according to $\operatorname{erfc}(x) \approx \frac{e^{-x^2}}{x\sqrt{\pi}}$.
			\item Consequently, the total error probability is dominated by the $\operatorname{erfc}(\cdot)$ term with the smallest argument, which corresponds to the worst-case destructive interference event ($1 - p/N$).
			\item For $p > 0$, since $1 - p/N < 1 + p/N$, the constructive term becomes negligible at high SNR, while destructive interference occurs in only $1$ out of $2$ equiprobable bit combinations of $U_2$ (multiplicity $1/2$).
		\end{itemize}
		
	\end{frame}
	
	% Slide 39b
	\begin{frame}{Exercise 2 - Part 3}
		\question{9. Derive an approximate performance curve based on $p$. (2/2)}
		
		\vspace{0.1cm}
		\begin{itemize}
			\item Accounting for the multiplicity factor $1/2$ (which multiplies the $1/2$ coefficient of $\operatorname{erfc}$), the high-SNR approximation for non-orthogonal codes ($p > 0$) is:
			\[
			P_{\text{approx}}(e) = \frac{1}{4}\operatorname{erfc}\!\left(\sqrt{\frac{E_b}{N_0}}\left(1 - \frac{p}{N}\right)\right)
			\]
			\item In the degenerate orthogonal case ($p = 0$), both terms in the exact expression share the identical argument $\sqrt{E_b/N_0}$.
			\item Both events contribute equally with multiplicity $2$ out of $2$, so the two $1/4$ terms sum to $1/2$, recovering the exact 2-PAM curve:
			\[
			P_{\text{approx}}(e) = \frac{1}{2}\operatorname{erfc}\!\left(\sqrt{\frac{E_b}{N_0}}\right)
			\]
		\end{itemize}
		
	\end{frame}
	
	% Slide 40
	\begin{frame}{Exercise 2 - Part 3}
		\question{10. Plot the analytical 2-PAM curve.}
		
		\begin{center}
			\includegraphics[width=0.72\textwidth, height=0.58\textheight, keepaspectratio]{A2_2_4_2PAM.png}
		\end{center}
	
		
	\end{frame}
	
	% Slide 41
	\begin{frame}{Exercise 2 - Part 3}
		\question{11. Plot the simulated curve obtained for that $p$.}
		
		\question{12. Add the analytic curve to the figure.}
		
		\question{13. Add the approximate curve to the figure.}
		
		\vspace{0.2cm}
		\begin{center}
			\includegraphics[width=0.72\textwidth, height=0.58\textheight, keepaspectratio]{A2_3_13_0.png}
		\end{center}
		
	\end{frame}
	
	% Slide 42
	\begin{frame}{Exercise 2 - Part 3}
		
		
		\vspace{0.2cm}
		\begin{center}
			\includegraphics[width=0.72\textwidth, height=0.58\textheight, keepaspectratio]{A2_3_13_2.png}
		\end{center}
		
	\end{frame}
	
	% Slide 43
	\begin{frame}{Exercise 2 - Part 3}
		
		
		\vspace{0.2cm}
		\begin{center}
			\includegraphics[width=0.72\textwidth, height=0.58\textheight, keepaspectratio]{A2_3_13_4.png}
		\end{center}
		
	\end{frame}
	
	% Slide 43
	\begin{frame}{Exercise 2 - Part 3}
		
		
		\vspace{0.2cm}
		\begin{center}
			\includegraphics[width=0.72\textwidth, height=0.58\textheight, keepaspectratio]{A2_3_13_12.png}
		\end{center}
		
	\end{frame}
	
	% Slide 44
	\begin{frame}{Exercise 2 - Part 3}
		\question{14. Comment on the result.}
		
		\vspace{0.1cm}
		\begin{itemize}
			\item Commenti ... 
		\end{itemize}
		
	\end{frame}
	
	% Slide 45
	\begin{frame}{Exercise 2 - Part 3}
		\question{15. Discuss the impact of $p$ on the performance.}
		
		\vspace{0.1cm}
		\begin{itemize}
			\item Commenti ... 
		\end{itemize}
		
	\end{frame}
	
	% Slide 46
	\begin{frame}[plain]
		\vfill
		\centering
		{\color{TealAccent}\rule{2cm}{2pt}}\\[0.4cm]
		{\LARGE\bfseries\color{NavyPrimary} Exercise 2 - Part 4 - CDMA Uplink}\\[0.2cm]
		\vfill
	\end{frame}
	
	% Slide 47
	\begin{frame}{Exercise 2 - Part 4}
		\question{16. What is the parameter p that determines the asymptotic performance of this scenario?}
		
		\vspace{0.1cm}
		\begin{itemize}
			\item Commenti ... 
		\end{itemize}
		
	\end{frame}
	
	% Slide 48
	\begin{frame}{Exercise 2 - Part 4}
		\question{17. Derive analytically the performance curve of this scenario with two users (and cyclic shift).}
		
		\vspace{0.1cm}
		\begin{itemize}
			\item Commenti ... 
		\end{itemize}
		
	\end{frame}
	
	% Slide 49
	\begin{frame}{Exercise 2 - Part 4}
		\question{18. Derive an approximate performance curve based on $p$.}
		
		\vspace{0.1cm}
		\begin{itemize}
			\item Commenti ... 
		\end{itemize}
		
	\end{frame}
	
	% Slide 50
	\begin{frame}{Exercise 2 - Part 4}
		\question{19. Plot the analytical 2-PAM curve.}
		
		\vspace{0.1cm}
		\begin{itemize}
			\item Grafico ... 
		\end{itemize}
		
	\end{frame}
	
	% Slide 51
	\begin{frame}{Exercise 2 - Part 4}
		\question{20. Plot the simulated curve obtained for that $p$.}
		
		\vspace{0.1cm}
		\begin{itemize}
			\item Grafico ... 
		\end{itemize}
		
	\end{frame}
	
	% Slide 52
	\begin{frame}{Exercise 2 - Part 4}
		\question{21. Add the analytic curve to the figure.}
		
		\vspace{0.1cm}
		\begin{itemize}
			\item Grafico ... 
		\end{itemize}
		
	\end{frame}
	
	% Slide 53
	\begin{frame}{Exercise 2 - Part 4}
		\question{22. Add the approximate curve to the figure.}
		
		\vspace{0.1cm}
		\begin{itemize}
			\item Grafico ... 
		\end{itemize}
		
	\end{frame}
	
	% Slide 54
	\begin{frame}{Exercise 2 - Part 4}
		\question{23. Comment on the result.}
		
		\vspace{0.1cm}
		\begin{itemize}
			\item Commenti ... 
		\end{itemize}
		
	\end{frame}
	
	% Slide 55
	\begin{frame}{Exercise 2}
		\question{24. Provide an overall discussion of the entire exercise.}
		
		\vspace{0.1cm}
		\begin{itemize}
			\item Commenti ... 
		\end{itemize}
		
	\end{frame}
	
	% Slide 56
	\begin{frame}[plain]
		\vfill
		\centering
		{\color{TealAccent}\rule{2cm}{2pt}}\\[0.4cm]
		{\LARGE\bfseries\color{NavyPrimary} Exercise 3 - Satellite Link Budget}\\[0.2cm]
		\vfill
	\end{frame}
	
	% Slide 57
	\begin{frame}{Exercise 3}
		\question{1. Justify your choices and report the sources of information used to select the system parameters..}
		
		\vspace{0.1cm}
		\begin{itemize}
			\item Commenti ... 
		\end{itemize}
		
	\end{frame}
	
	% Slide 58
	\begin{frame}[plain]
		\vfill
		\centering
		{\color{TealAccent}\rule{2cm}{2pt}}\\[0.4cm]
		{\LARGE\bfseries\color{NavyPrimary} Exercise 3 - Part 1 - Satellite-to-User Distance}\\[0.2cm]
		\vfill
	\end{frame}
	
	% Slide 59
	\begin{frame}{Exercise 3 - Part 1}
		\question{2. Compute the distance $d$ for the selected elevation angle $\theta_e$.}
		
		\vspace{0.1cm}
		\begin{itemize}
			\item Commenti ... 
		\end{itemize}
		
	\end{frame}
	
	% Slide 60
	\begin{frame}{Exercise 3 - Part 1}
		\question{3. Plot the distance $d$ and the free-space path loss in dB versus the elevation angle $\theta_e$ for $10^\circ \leq \theta_e \leq 90^\circ$.}
		
		\vspace{0.1cm}
		\begin{itemize}
			\item Commenti ... 
		\end{itemize}
		
	\end{frame}
	
	% Slide 61
	\begin{frame}{Exercise 3 - Part 1}
		\question{4. Comment on the results.}
		
		\vspace{0.1cm}
		\begin{itemize}
			\item Commenti ... 
		\end{itemize}
		
	\end{frame}
	
	% Slide 62
	\begin{frame}[plain]
		\vfill
		\centering
		{\color{TealAccent}\rule{2cm}{2pt}}\\[0.4cm]
		{\LARGE\bfseries\color{NavyPrimary} Exercise 3 - Part 2 - Link Budget}\\[0.2cm]
		\vfill
	\end{frame}
	
	% Slide 63
	\begin{frame}{Exercise 3 - Part 2}
		\question{5. Compute the TX and RX antenna gains $G_{TX}$ and $G_{RX}$ (in dBi) using realistic values for their efficiencies.}
		
		\vspace{0.1cm}
		\begin{itemize}
			\item Commenti ... 
		\end{itemize}
		
	\end{frame}
	
	% Slide 64
	\begin{frame}{Exercise 3 - Part 2}
		\question{6. Compute the EIRP (in dBW).}
		
		\vspace{0.1cm}
		\begin{itemize}
			\item Commenti ... 
		\end{itemize}
		
	\end{frame}
	
	% Slide 65
	\begin{frame}{Exercise 3 - Part 2}
		\question{7. Compute the free-space path loss (in dB).}
		
		\vspace{0.1cm}
		\begin{itemize}
			\item Commenti ... 
		\end{itemize}
		
	\end{frame}
	
	% Slide 66
	\begin{frame}{Exercise 3 - Part 2}
		\question{8. Consider additional losses $L_{add} = 3$ dB (accounting for atmospheric, pointing and polarization losses) and compute the total losses (in dB).}
		
		\vspace{0.1cm}
		\begin{itemize}
			\item Commenti ... 
		\end{itemize}
		
	\end{frame}
	
	% Slide 67
	\begin{frame}{Exercise 3 - Part 2}
		\question{9. Compute the received power $P_{RX}$ (in dBW).}
		
		\vspace{0.1cm}
		\begin{itemize}
			\item Commenti ... 
		\end{itemize}
		
	\end{frame}
	
	% Slide 68
	\begin{frame}{Exercise 3 - Part 2}
		\question{10. Compute the noise power $P_N$ (in dBW).}
		
		\vspace{0.1cm}
		\begin{itemize}
			\item Commenti ... 
		\end{itemize}
		
	\end{frame}
	
	% Slide 69
	\begin{frame}{Exercise 3 - Part 2}
		\question{11. Compute the signal-to-noise ratio SNR (in dB).}
		
		\vspace{0.1cm}
		\begin{itemize}
			\item Commenti ... 
		\end{itemize}
		
	\end{frame}
	
	% Slide 70
	\begin{frame}{Exercise 3 - Part 2}
		\question{12. Summarize all the numerical results in a link budget table.}
		
		\vspace{0.1cm}
		\begin{itemize}
			\item Table ... 
		\end{itemize}
		
	\end{frame}
	
	% Slide 71
	\begin{frame}[plain]
		\vfill
		\centering
		{\color{TealAccent}\rule{2cm}{2pt}}\\[0.4cm]
		{\LARGE\bfseries\color{NavyPrimary} Exercise 3 - Part 3 - Bit Rate}\\[0.2cm]
		\vfill
	\end{frame}
	
	
	% Slide 72
	\begin{frame}{Exercise 3 - Part 3}
		\question{13. Select the MODCOD providing the highest spectral efficiency among those reported in the table, compatible with the available $E_s/N_0$.}
		
		\vspace{0.1cm}
		\begin{itemize}
			\item Commenti ... 
		\end{itemize}
		
	\end{frame}
	
	% Slide 73
	\begin{frame}{Exercise 3 - Part 3}
		\question{14. Report its spectral efficiency and compute the corresponding bit rate $R_b$.}
		
		\vspace{0.1cm}
		\begin{itemize}
			\item Commenti ... 
		\end{itemize}
		
	\end{frame}
	
	% Slide 74
	\begin{frame}{Exercise 3 - Part 3}
		\question{15. Assuming an average bit rate of $5$ Mbit/s per HD television channel, estimate how many television channels can be simultaneously broadcast by the considered satellite carrier, neglecting additional service overhead.}
		
		\vspace{0.1cm}
		\begin{itemize}
			\item Commenti ... 
		\end{itemize}
		
	\end{frame}
	
	% Slide 75
	\begin{frame}{Exercise 3 - Part 3}
		\question{16. Compute the Shannon capacity $C = B\log_2(1+SNR)$.}
		
		\vspace{0.1cm}
		\begin{itemize}
			\item Commenti ... 
		\end{itemize}
		
	\end{frame}
	
	% Slide 76
	\begin{frame}{Exercise 3 - Part 3}
		\question{17. Compare $R_b$ with $C$ and comment on the results.}
		
		\vspace{0.1cm}
		\begin{itemize}
			\item Commenti ... 
		\end{itemize}
		
	\end{frame}
	
	% Slide 77
	\begin{frame}[plain]
		\vfill
		\centering
		{\color{TealAccent}\rule{2cm}{2pt}}\\[0.4cm]
		{\LARGE\bfseries\color{NavyPrimary} Exercise 3 - Part 4 - System Trade-offs}\\[0.2cm]
		\vfill
	\end{frame}
	
	% Slide 78
	\begin{frame}{Exercise 3 - Part 4}
		\question{18. Select and justify a suitable range of variation for the power and the bandwidth.}
		
		\vspace{0.1cm}
		\begin{itemize}
			\item Commenti ... 
		\end{itemize}
		
	\end{frame}
	
	% Slide 79
	\begin{frame}{Exercise 3 - Part 4}
		\question{19. Plot the bit rate versus the transmitted power $P_{TX}$.}
		
		\vspace{0.1cm}
		\begin{itemize}
			\item Plot ... 
		\end{itemize}
		
	\end{frame}
	
	% Slide 80
	\begin{frame}{Exercise 3 - Part 4}
		\question{20. Plot the bit rate versus the bandwidth $B$.}
		
		\vspace{0.1cm}
		\begin{itemize}
			\item Plot ... 
		\end{itemize}
		
	\end{frame}
	
	% Slide 81
	\begin{frame}{Exercise 3 - Part 4}
		\question{21. Discuss the results.}
		
		\vspace{0.1cm}
		\begin{itemize}
			\item Commenti ... 
		\end{itemize}
		
	\end{frame}
	
	
	% Slide finale
	\begin{frame}[plain]
		\vfill
		\centering
		{\color{TealAccent}\rule{2.5cm}{2pt}}\\[0.5cm]
		{\Huge\bfseries\color{NavyPrimary} Grazie!}\\[0.3cm]
		\vfill
	\end{frame}
	
\end{document}